% !PRINTED 300
\weberpage{300}

Let $\alpha$ and $\beta > \alpha$ again be two real numbers, which determine the interval $(\alpha, \beta)$ in which the number of real roots of $f(x)$ is to be counted.

We shall first assume that in the interval $(\alpha, \beta)$ no two terms of the sequence (1) vanish simultaneously, and that for $x = \alpha$, $x = \beta$ no term of the sequence vanishes.

If we now let $x$ increase continuously from $\alpha$ to $\beta$, a change can occur in the succession of signs of the sequence (1) only when one of the functions passes through zero. When $f^{(\nu)}(x)$ passes through zero for $x = \xi$, then $f^{(\nu)}(x)$ passes from negative to positive or from positive to negative values, according as $f^{(\nu+1)}(x)$ is positive or negative, and thus a change of sign between $f^{(\nu)}$ and $f^{(\nu+1)}$ is lost in the passage through $\xi$. This holds also when $f^{(\nu)}$ is the function $f(x)$ itself. But when $f^{(\nu)}$ is one of the derivatives, it is preceded by a function $f^{(\nu-1)}$, and since $f^{(\nu-1)}$ is by hypothesis not zero for $x = \xi$, and $f^{(\nu)}$ changes its sign in the passage through $\xi$, either a loss or a gain of one change of sign takes place between $f^{(\nu-1)}$ and $f^{(\nu)}$ in the passage through $\xi$. Thus when an inner term of the sequence (1) passes through zero, the number of changes of sign remains unaltered, or two changes of sign are lost.

But if $f(x)$ itself passes through zero, one change of sign is lost.

From this follows the theorem:

I. \gesperrt{The number of roots of $f(x)$ lying between $\alpha$ and $\beta$ is at most as great as the number of changes of sign lost between $\alpha$ and $\beta$, and if it is smaller, the difference is an even number.}

Making use of a formula, we can also say:

If $V(x)$ is the number of changes of sign (variations) that the sequence (1) presents for any value $x$, then the number of roots of $f(x)$ lying between $\alpha$ and $\beta$ is

\weq{2}{V(\alpha) - V(\beta) - 2h,}

where $h$ is a non-negative integer.

The proof of theorem I. still requires a supplement for the case in which several successive terms of the sequence (1) vanish simultaneously.
